\documentclass[12pt]{article}
\usepackage{amsfonts,amsthm,amsmath,amssymb,mathtools}
\usepackage{mathrsfs}
\usepackage{enumitem}
\usepackage{tikz-cd}
\usepackage{geometry}
\usepackage{graphicx}
\usepackage{mlmodern}

\geometry{letterpaper, portrait, margin=1in}

\setcounter{section}{0}

\renewcommand{\thesection}{\S\arabic{section}}
\renewcommand{\thesubsection}{\arabic{section}}
\DeclareMathOperator{\im}{im}
\DeclareMathOperator{\id}{id}
\DeclareMathOperator{\ad}{ad}
\DeclareMathOperator{\tr}{tr}
\newcommand{\gl}{\mathfrak{gl}}
\newcommand{\so}{\mathfrak{so}}
\renewcommand{\sl}{\mathfrak{sl}}
\renewcommand{\th}{^\mathrm{th}}

\newtheorem{thm}{Theorem}
\newenvironment{theorem}[1]{
	\renewcommand{\thethm}{#1}
	\begin{thm}
		}{\end{thm}}

\newtheorem{lma}{Lemma}
\newenvironment{lemma}[1]{
	\renewcommand{\thelma}{#1}
	\begin{lma}
		}{\end{lma}}

\theoremstyle{definition}
\newtheorem{ex}{}
\newenvironment{exercise}[1]{
	\renewcommand{\theex}{#1}
	\begin{ex}
		}{\end{ex}}

\title{Lie Theory: Chapter 4 Homework}
\author{Connor Mooneyhan}
\date{}

\begin{document}

\maketitle

\begin{ex}
	Let $L_1$ and $L_2$ be Lie algebras, and let $\varphi : L_1 \to L_2$ be a surjective Lie algebra homomorphism.
	True or false:\begin{enumerate}[label=(\alph*)]
		\item $\varphi(L_1') = L_2'$
		\item $\varphi(Z(L_1)) = Z(L_2)$
	\end{enumerate}
\end{ex}
\begin{proof}\
	\begin{enumerate}[label=(\alph*)]
		\item Yes.
		      We have \begin{align*}
			      \varphi\left(\sum a_i [x_i, y_i]\right) = \sum a_i \varphi([x_i, y_i]) = \sum a_i [\varphi(x_i), \varphi(y_i)] \in L_2',
		      \end{align*}
		      so $\varphi(L_1') \subseteq L_2'$.

		      On the other hand, for any $x, y \in L_2$, we have $x', y' \in L_1$ such that $\varphi(x') = x$ and $\varphi(y') = y$, so \begin{align*}
			      \sum a_i [x_i, y_i] & = \sum a_i [\varphi(x_i'), \varphi(y_i')]                        \\
			                          & = \sum a_i \varphi([x_i', y_i'])                                 \\
			                          & = \varphi\left( \sum a_i [x_i', y_i'] \right) \in \varphi(L_1'),
		      \end{align*}
		      so $L_2' \subseteq \varphi(L_1')$, so $\varphi(L_1') = L_2'$.
		\item If $x \in Z(L_1)$, then given $y \in L_2$, we have $y = \varphi(y')$ for some $y' \in L_1$, so \begin{align*}
			      [\varphi(x), y] & = [\varphi(x), \varphi(y')] \\
			                      & = \varphi([x, y'])          \\
			                      & = \varphi(0)                \\
			                      & = 0,
		      \end{align*}
		      so we have $\varphi(x) \in Z(L_2)$, and thus $\varphi(Z(L_1)) \subseteq Z(L_2)$.

		      The other direction does not hold, however.
		      Consider the two-dimensional nonabelian Lie algebra $L_1 = \mathrm{span}\lbrace x, y \rbrace$ with $[x, y] = x$.
		      Consider the quotient map $\varphi : L_1 \to L_1 / \mathrm{span}\lbrace x \rbrace$, so $L_2 = L_1 / \mathrm{span}\lbrace x\rbrace$.
		      Then $Z(L_1) = 0$, so $\varphi(Z(L_1)) = 0$, but $L_2$ is abelian since it is one-dimensional, so $Z(L_2) = L_2$.
		      Thus $Z(L_2) \not\subseteq \varphi(Z(L_1))$ in general.
	\end{enumerate}
\end{proof}

\begin{ex}
	Use Lemma 4.4 to show that if $L$ is a Lie algebra then $L$ is solvable if and only if $\ad L$ is a solvable subalgebra of $\gl(L)$.
	Show that this result also holds if we replace ``solvable'' with ``nilpotent.''
\end{ex}
\begin{proof}
	Note that since $\ad : L \to \gl(L)$ has kernel $Z(L)$, $\ad L \cong L / Z(L)$. Lemma 4.4(a) then tells us that as a homomorphic image of $L$, $\ad L$ is solvable if $L$ is solvable.
	We then see that since $Z(L)$ is itself a solvable ideal since $Z(L)^2$, if $\ad L \cong L/Z(L)$ is solvable, Lemma 4.4(b) tells us that $L$ is solvable.

	For the nilpotent versions, note that since $\ad : L \to \gl(L)$ is a homomorphism, we have $(\ad L)^n = \ad (L^n )$ for any $n$.
	Thus if $L$ is nilpotent with $L^m = 0$, we have \begin{equation*}
		(\ad L)^m = \ad(L^m) = \ad 0 = 0,
	\end{equation*}
	and if instead $\ad L$ is nilpotent with $(\ad L)^m = 0$, we have $\ad (L^m) = 0$, implying that $L^m \subseteq Z(L)$, so $L^{m+1} = 0$ and $L$ is nilpotent.
\end{proof}

\begin{ex}
	Let $L = \mathbf{n}(n, F)$.
	Show that $L^k$ has a basis consisting of all the matrix units $e_{ij}$ with $j - i > k$.
	Hence show that $L$ is nilpotent.
	What is the smallest $m$ such that $L^m = 0$?
\end{ex}
\begin{proof}
	As we have seen, when we bracket two elements of $L$ together, the bottommost diagonal line of entries with nonzero elements becomes all-zero.
	Thus $L^1$ consists of matrices which are strictly upper-triangular \textit{and} have zeroes in entries $i,j$ where $j - i > 1$.
	Similarly, $L^k$ has zeroes everywhere except in entries $i,j$ where $j - i > k$.
	Thus it has a basis of such entries.
	If $\dim L = n$, then $L^n$ would need to have a basis of matrix elements $e_{ij}$ where $j-i > n$, but this is impossible, so $L^n = 0$ and $L$ is nilpotent.
	The lowest such $m$ is in fact $m=n - 1$, since we need the top-right entry at $1, n$ to be zero.
\end{proof}

\begin{ex}
	Let $L = \mathbf{b}(n, F)$ be the Lie algebra of upper triangular $n \times n$ matrices over a field $F$.
	\begin{enumerate}[label=(\roman*)]
		\item Show that $L' = \mathbf{n}(n, F)$.
		\item More generally, show that $L^{(k)}$ has a basis consisting of all the matrix units $e_{ij}$ with $j - i \geq 2^{k-1}$.
		      (The commutator formula for the $e_{ij}$ given in $\S 1.2$ will be helpful.)
		\item Hence show that $L$ is solvable.
		      What is the smallest $m$ such that $L^{(m)} = 0$?
		\item Show that if $n \geq 2$ then $L$ is not nilpotent.
	\end{enumerate}
\end{ex}
\begin{proof}\
	\begin{enumerate}[label=(\roman*)]
		\item When matrix-multiplying two upper triangular matrices, the diagonal elements multiply together (i.e. $(AB)_{k,k} = A_{k, k}B_{k,k}$).
		      Thus when we take the bracket, we get \begin{align*}
			      [A, B]_{k,k} & = (AB - BA)_{k,k}                 \\
			                   & = (AB)_{k,k} - (BA)_{k,k}         \\
			                   & = A_{k,k}B_{k,k} - B_{k,k}A_{k,k} \\
			                   & = 0,
		      \end{align*}
		      so there are zeroes along the diagonal.
		      We can still get nonzero entries above the diagonal though, so indeed $L' = \mathbf{n}(n, F)$.
		\item The aforementioned commutator formula tells us that $[e_{ij}, e_{kl}] = \delta_{jk}e_{il} - \delta_{il}e_{kj}$.
		      If $j = k$ and $i = l$ (so that the deltas are valued at 1) and $i = k$ and $l = j$ (so that $e_{il} = e_{kj}$), then $[e_{ij}, e_{kl}] = 0$, but in this case $i = j = k = l$, so we are just talking about the bracket of an element with itself.
		      The other way $[e_{ij}, e_{kl}] = 0$ is if $j \neq k$ and $i \neq l$ so that both deltas are zero.
	\end{enumerate}
\end{proof}

\begin{ex}
	Show that a Lie algebra is semisimple if and only if it has no non-zero abelian ideals.
\end{ex}
\begin{proof}
	Let $L$ be a Lie algebra.
	First, since every abelian Lie algebra is solvable, every abelian ideal is solvable, so if there are nonzero abelian ideals of $L$, then $L$ is not semisimple.
	Now if $L$ is not semisimple, then let $I$ be a nonzero solvable ideal of $L$ and let $m$ be the smallest integer such that $I^{(m)} = 0$.
	Then $I^{(m-1)} \neq 0$ and $(I^{(m-1)})' = 0$, so $I^{(m-1)}$ is abelian, but it is also an ideal and nonzero.
\end{proof}

\setcounter{ex}{6}

\begin{ex}
	Let $L_1$ and $L_2$ be Lie algebras, and let $\varphi : L_1 \to L_2$ be a surjective Lie algebra homomorphism.
	Show that if $h \in L_1$ and $\ad h$ is diagonalizable then $\ad \varphi(h)$ is diagonalizable.
\end{ex}
\begin{proof}
	Let $x_1, \dots, x_n$ be a basis for $L_1$ consisting of eigenvectors of $\ad h$ and let $\lambda_1, \dots, \lambda_n$ be their corresponding eigenvalues, so that $[h, x_i] = \lambda x_i$ for each $i$.
	Then for each $i$, \begin{align*}
		[\varphi(h), \varphi(x_i)] & = \varphi([h, x_i])      \\
		                           & = \varphi(\lambda_i x_i) \\
		                           & = \lambda_i\varphi(x_i),
	\end{align*}
	so $\varphi(x_i) = 0$ or $\varphi(x_i)$ is an eigenvector of $\ad \varphi(h)$.
	We know that $\left\lbrace\varphi(x_1), \dots, \varphi(x_n)\right\rbrace$ does span $L_2$ since $\left\lbrace x_1, \dots, x_n\right\rbrace$ spans $L_1$ and $\varphi$ is surjective, so it must contain a basis for $L_2$, in particular a basis of eigenvectors, and thus $\ad \varphi(h)$ is diagonalizable.
\end{proof}

\begin{ex}
	Prove directly that $\sl(n, \mathbb{C})$ is a simple Lie algebra for $n \geq 2$.
\end{ex}
\begin{proof}
	Let $I$ be a nonzero ideal of $\sl(n, \mathbb{C})$, and let $x \in I$ such that $x \neq 0$.
	If $x$ is diagonal, then the trace is zero so not all of its diagonal entries can be equal, so let $i \neq j$ such that $x_{ii} \neq x_{jj}$ and both are nonzero.
	Then \begin{align*}
		[x, e_{ij}] & = xe_{ij} - e_{ij}x            \\
		            & = x_{ii}e_{ij} - e_{ij}x_{jj}  \\
		            & = (x_{ii}-x_{jj})e_{ij} \in I,
	\end{align*}
	so since a multiple of $e_{ij}$ is in $I$, $e_{ij}$ itself is in $I$.
	If $x$ is not diagonal, then there exists a nonzero entry $x_{ji}$ where $j \neq i$, and we see that \begin{align*}
		[[x, e_{ij}], e_{ij}] & = [x, e_{ij}]e_{ij} - e_{ij}[x, e_{ij}]                         \\
		                      & = (xe_{ij} - e_{ij}x)e_{ij} - e_{ij}(xe_{ij} - e_{ij}x)         \\
		                      & = xe_{ij}e_{ij} - e_{ij}xe_{ij} - e_{ij}xe_{ij} + e_{ij}e_{ij}x \\
		                      & = -2e_{ij}xe_{ij} \hspace{1em} \text{(since $e_{ij}^2 = 0$)}    \\
		                      & = -2x_{ji}e_{ij} \in I,
	\end{align*}
	so we again have $e_{ij} \in I$.
	Either way, we have established $i \neq j$ such that $e_{ij} \in I$.

	Now we see that $[e_{ij}, e_{kl}] \in I$ for any $k$ and $l$, so using the commutator formula \begin{equation*}
		[e_{ij}, e_{kl}] = \delta_{jk}e_{il} - \delta_{li}e_{kj},
	\end{equation*}
	so $[e_{ij}, e_{jl}] = e_{il} \in I$ for any $l$ so long as $l \neq i$, and $[e_{ij}, e_{ki}] = -e_{kj} \in I \implies e_{kj} \in I$ for any $k$ where $k \neq j$.
	From there, we use similar tricks to show that every off-diagonal matrix unit $e_{kl}$ is in $I$.
	Finally, now that we see that $e_{m,m+1} \in I$ whenever $m < n$, we have $[e_{m, m+1}, e_{m+1, m}] = e_{m,m} - e_{m+1,m+1} \in I$.
	We have now shown that every standard basis element of $\sl(n, \mathbb{C})$ is in $I$, so $I = \sl(n, \mathbb{C})$, proving that $\sl(n, \mathbb{C})$ is simple.
\end{proof}

\begin{ex}\
	\begin{enumerate}[label=(\roman*)]
		\item Show that $\det(I + \varepsilon X)$ is a polynomial in $\varepsilon$ of degree at most $n$ with the first two terms \begin{equation*}
			      \det(I + \varepsilon X) = 1 + (\tr X)\varepsilon + \cdots.
		      \end{equation*}
		      If we neglect all powers of $\varepsilon$ except 1 and $\varepsilon$, then we obtain the statement \begin{equation*}
			      I + \varepsilon X \in \mathrm{SL}(n, \mathbb{C}) \iff X \in \sl(n, \mathbb{C}).
		      \end{equation*}
		\item \begin{enumerate}[label=(\alph*)]
			      \item Let $S$ be an $n \times n$ matrix.
			            Let $(-,-)$ denote the complex bilinear form with matrix $S$.
			            Show that if we let $G_S(n, \mathbb{C})$ be the set of invertible matrices $A$ such that $(Au, Av) = (u, v)$ for all $u, v \in \mathbb{C}^n$, then $G_S(n, \mathbb{C})$ is a group.
			      \item Show that if we perform the construction in part (i) with $G_S(n, \mathbb{C})$ in place of $\mathrm{SL}(n, \mathbb{C})$, we obtain $\gl_S(n, \mathbb{C})$.
		      \end{enumerate}
		\item \begin{enumerate}[label=(\alph*)]
			      \item An invertible matrix $A$ is customarily said to be \textit{orthogonal} if $A^tA = AA^t = I$; that is, if $A^{-1} = A^t$. Show that the set of $n\times n$ orthogonal matrices with coefficients in $\mathbb{C}$ is the group $G_I(n, \mathbb{C})$ and that the associated Lie algebra, $\gl_I(n, \mathbb{C})$, is the space of all anti-symmetric matrices.
			      \item Prove that $\gl_I(n, \mathbb{C}) \cong \so(n, \mathbb{C})$.
			            \textit{Hint}: Use Exercise 2.11.
		      \end{enumerate}
		\item A bilinear form on a vector space $V$ is said to be \textit{symplectic} if $(v, v) = 0$ for all $v \in V$.
		      Show that \begin{equation*}
			      S= \begin{pmatrix}0 & I_\ell \\ -I_\ell & 0\end{pmatrix}
		      \end{equation*}
		      is the matrix of a non-degenerate symplectic bilinear form on a $2\ell$-dimensional space.
		      The associated Lie algebra is $\gl_S(2\ell, \mathbb{C}) = \mathfrak{sp}(2\ell, \mathbb{C})$.
	\end{enumerate}
\end{ex}
\begin{proof}\
	\begin{enumerate}[label=(\roman*)]
		\item If $\lambda$ is an eigenvalue of $X$ and $x$ is an eigenvector corresponding to $\lambda$, then \begin{equation*}
			      (I+\varepsilon X)x = x + \varepsilon\lambda x = (1+\varepsilon\lambda)x,
		      \end{equation*} so $1 + \varepsilon\lambda$ is an eigenvalue of $I + \varepsilon X$.
		      Thus if $\lambda_1, \dots, \lambda_n$ are the eigenvalues of $X$ repeated according to multiplicity, then \begin{align*}
			      \det(I + \varepsilon X) & = \prod_{1 \leq i \leq n} (1 + \varepsilon\lambda_i),
		      \end{align*}
		      which is indeed a polynomial of degree $n$ in $\varepsilon$ whose first term is $1$ and whose second term is \begin{equation*}
			      \sum_{i}\varepsilon\lambda_i = (\tr X) \varepsilon.
		      \end{equation*}
		% \item \begin{enumerate}[label=(\alph*)]
		% 	      \item .
		% 	      \item .
		%       \end{enumerate}
		% \item \begin{enumerate}[label=(\alph*)]
		% 	      \item .
		% 	      \item .
		%       \end{enumerate}
		% \item .
	\end{enumerate}
\end{proof}

\end{document}
